% - In the regression setting of the original CSA paper, envelopes are typically convex because the score space is 
% unbounded and the data does not cluster into separated groups\\
% - In our classification setting, the score space is bounded in $[0,1]^K$ and — crucially — true NC vectors for different 
% classes occupy distinct, well-separated regions of this space (example)\\
% - In the classification, 
% % non-convex envelopes are not only permitted but actively beneficial 
% produce tighter prediction sets without sacrificing the coverage guarantee\\
% - The guarantee still holds because it depends only on the exchangeability of the data and the quantile construction, 
% not on the shape of the envelope
% % This is the central geometric motivation for the radial and strip methods developed in section 3.5
% Briefly contrast with collapsed 1D: if the data is genuinely one-dimensional (all dimensions carry the same signal), 
%collapsing to a scalar is sufficient and non-convexity gives no additional benefit


\subsection{Non-Convex Envelopes in Classification}
\label{sec3.2: non-convex envelopes}
\noindent Section~\ref{sec2.4.3: csa limitations} identified convexity of the envelope $E$ as one of the main limitations of the 
CSA method when applied to classification. In regression, the true target values usually form a single, continuous region. A 
convex shape is therefore a natural fit because it fills in the region without creating artificial gaps. In classification, 
however, score vectors for different classes tend to cluster in separate groups. So forcing a single convex shape around these 
separate clusters stretches the region over empty space, leading to unnecessarily large prediction sets without improving 
coverage. Here, we will describe the alternate method that we have used in this report.

%%%%%%%%%%%%%%%%%%%%%%%%%%%%%%%%%%

\vspace{0.15cm}
\noindent Let $A, B \subset [0,1]^K$ denote the regions occupied by the transformed nonconformity vectors 
$\phi_h(\mathbf{s})$ of two distinct true labels, and suppose $A$ and $B$ are separated, i.e. there exists 
$\varepsilon > 0$ such that $\|\mathbf{a} - \mathbf{b}\| \ge \varepsilon$ for all $\mathbf{a} \in A, \mathbf{b} \in B$. 
If $E$ is required to be convex and to contain both $A$ and $B$, then necessarily $E \supseteq \operatorname{conv}(A 
\cup B)$. Since $A$ and $B$ are separated,
\begin{equation*}
\operatorname{vol}\big(\operatorname{conv}(A \cup B) \setminus (A \cup B)\big) > 0,
\end{equation*}
so any convex envelope containing both classes strictly contains a region of positive volume belonging to neither 
class. Because membership in the scaled envelope $\hat{t} \cdot E$ determines set inclusion 
(Equation~\eqref{eq: prediction set definition}), this excess volume directly inflates $|C(\mathbf{s})|$ for any test 
point whose transformed score lands in it, without improving coverage. Therefore, convexity is not just an unnecessary 
restriction here but it also affects efficiency whenever the true classes occupy well-separated regions of score space, which 
the empirical score statistics of \todo{refer }confirm is the case for gram-type classification (gram-neg and gram-pos scores 
concentrate near $0.07$ and $0.65$ respectively).

After applying the label-dependent transformation $\phi_y$ (Section~\ref{sec3.3: nc transformation}), the nonconformity score 
vectors of true infectors concentrate near the origin in $\mathbb{R}^K_+$. Crucially, the clusters for different classes are not 
intermingled. In our gram-type task, gram-negative phages have mean raw scores $\approx 0.06$--$0.08$ while gram-positive phages 
have mean raw scores $\approx 0.63$--$0.81$. Under the transformation \eqref{eq: gram-nc-transform}, the gram-negative NC 
vectors live in $[0,0.1]^K$ whereas the gram-positive NC vectors live in $[0.2,0.4]^K$. The two populations are separated by a 
gap of width $\approx 0.1$--$0.2$ along every dimension.

Formally, let $\mathcal{C}_y \subseteq \mathbb{R}^K_+$ denote the support of $\phi_y(\mathbf{S})$ when $Y=y$ is the true label. In our setting,
\begin{equation}
\label{eq: separated supports}
\mathcal{C}_{\text{neg}} \cap \mathcal{C}_{\text{pos}} = \emptyset,
\qquad
\operatorname{dist}(\mathcal{C}_{\text{neg}}, \mathcal{C}_{\text{pos}}) > 0.
\end{equation}
If we construct a convex envelope $E_{\text{conv}} = \operatorname{conv}(\mathcal{C}_{\text{neg}} \cup \mathcal{C}_{\text{pos}})$, then by Carathéodory's theorem every point of $E_{\text{conv}}$ is a convex combination of at most $K+1$ points from $\mathcal{C}_{\text{neg}} \cup \mathcal{C}_{\text{pos}}$. Because the supports are separated, $E_{\text{conv}}$ contains points that are not in $\mathcal{C}_{\text{neg}} \cup \mathcal{C}_{\text{pos}}$---the ``empty region'' between the clusters. At test time, any NC vector falling in this gap is included in the prediction set for \emph{both} classes, producing a size-2 set even though no true infector ever occupies that region.




%%%%%%%%%%%%%%%%%%%%%%%%%%%%%%%%%%%%%%%%%%%%%

\vspace{0.15cm}
\noindent Thus, we drop the convexity requirement on $E$, and only keep the properties needed for the coverage 
guarantee of Theorem~\ref{thm: marginal coverage} to hold: the guarantee depends only on the exchangeability of 
calibration and test points and on the quantile construction of $\hat{t}$ \cite{ochoarivera2024conformal}. Any measurable 
region $E$ for which membership can be evaluated at test time yields a valid coverage guarantee once $\hat{t}$ is calibrated on 
it.

\vspace{0.15cm}
\noindent One structural property is worth isolating, since one of the three envelope geometries of 
Section~\ref{sec3.4: envelope construction} satisfies it. A region $E$ is \emph{star-shaped} with respect to the 
origin if
\begin{equation*}
\mathbf{x} \in E \implies \lambda \mathbf{x} \in E, \qquad \forall \lambda \in [0,1].
\end{equation*}
Star-shapedness is strictly weaker than convexity, but it is exactly the property needed for the scalar 
$\tau(\mathbf{s})$ of Section~\ref{sec2.4.1: ensemble score vectors} to be well-defined as a single scaling threshold: 
if $\mathbf{s} \in t \cdot E$ for some $t > 0$, star-shapedness guarantees $\mathbf{s} \in t' \cdot E$ for every 
$t' \ge t$, so the set $\{t > 0 : \mathbf{s} \in t \cdot E\}$ is an interval and its infimum, $\tau(\mathbf{s})$, is a 
single well-defined boundary rather than one of possibly several disconnected thresholds.

\vspace{0.15cm}
\noindent The Radial envelope of Section~\ref{sec3.4: envelope construction} is star-shaped by construction: each 
directional boundary is a radius $\tilde{q}_m$ along a fixed direction $\mathbf{u}_m$ from the origin, so a point 
$\mathbf{s}$ satisfying $s_k \le \tilde q_m u_{mk}$ for all $k$, for some sector $m$, trivially satisfies 
$\lambda s_k \le \tilde q_m u_{mk}$ for any $\lambda \in [0,1]$ — downscaling toward the origin preserves the 
inequality coordinate-wise. It is not convex, since $\tilde q_m$ varies across directions, but this non-convexity is a 
deliberate consequence of letting the envelope follow the true anisotropic shape of the calibration data, not a 
defect. The Strip envelope, by contrast, conditions each dimension's boundary on the bin index of the point's own 
coordinate; its monotonicity propagation (Section~\ref{sec3.4: envelope construction}) enforces a weaker, one-sided 
ordering across bins and does not in general guarantee star-shapedness. We therefore describe Strip only as 
non-convex and data-adaptive, without claiming the stronger property for it.







\paragraph{Efficiency loss from convexity.} Let $\operatorname{Vol}(\cdot)$ denote the Lebesgue measure in $\mathbb{R}^K$. The volume of the convex envelope satisfies
\begin{equation}
\label{eq: volume inequality}
\operatorname{Vol}(E_{\text{conv}})
\;\geq\;
\operatorname{Vol}(\mathcal{C}_{\text{neg}}) + \operatorname{Vol}(\mathcal{C}_{\text{pos}})
+ \operatorname{Vol}(\text{gap}),
\end{equation}
where the gap term is strictly positive under \eqref{eq: separated supports}. Because the CSA scaling factor $\hat{t}$ is chosen to achieve coverage $1-\alpha$ on $S_2$, a larger envelope volume directly translates into larger expected prediction sets $\mathbb{E}[|C(\mathbf{s})|]$. The efficiency loss is most severe when the clusters are far apart relative to their individual diameters---precisely the regime we observe in phage-host score data.

\paragraph{Non-convex envelopes preserve coverage.} Dropping convexity does not invalidate the conformal guarantee. Theorem~\ref{thm: marginal coverage} holds for \emph{any} measurable envelope $E \subseteq \mathbb{R}^K_+$, provided the scaling factor $\hat{t}$ is computed as the $(1-\alpha)$ quantile of the $\tau$-values over $S_2$. The proof relies only on exchangeability and the rank argument of Lemma~\ref{lem: uniform rank}; it places no restriction on the geometric shape of $E$. We may therefore design envelopes that adapt to the local geometry of each class cluster without sacrificing the finite-sample coverage guarantee.

Figure~\ref{fig: nonconvex} illustrates this principle in two dimensions. Panel (a) shows a convex hull enclosing both gram-negative and gram-positive NC clusters; the shaded red region represents the wasted coverage between the clusters. Panel (b) shows a non-convex envelope that wraps tightly around each cluster, eliminating the gap and reducing the expected prediction-set size.

% \begin{figure}[htbp]
% \centering
% \includegraphics[width=0.95\textwidth]{figure_nonconvex_envelopes.png}
% \caption{Envelope geometry in nonconformity score space. \textbf{(a)} A convex envelope (black dashed boundary) enclosing both gram-negative (blue) and gram-positive (orange) NC clusters. The red shaded region lies inside the convex hull but contains no data points from either class; any test point falling here receives a size-2 prediction set. \textbf{(b)} A non-convex envelope consisting of two separate regions fitted to each cluster. The same coverage guarantee holds, but the expected prediction-set size is strictly smaller.}
% \label{fig: nonconvex}
% \end{figure}

\paragraph{Practical construction.} In Sections~\ref{sec3.4.1: radial}--\ref{sec3.4.3: collapsed} we present three envelope geometries. The Radial and Strip methods are explicitly non-convex: the Radial envelope assigns different boundary radii to different angular sectors, while the Strip envelope allows the admissible range of each dimension to vary conditionally on the value of another dimension. The Collapsed 1D method is a scalar baseline that discards geometric structure entirely. All three methods satisfy the coverage guarantee \eqref{eq: coverage-guarantee-31} by construction; their efficiency differs according to how well they approximate the true class-conditional supports $\mathcal{C}_y$.